% !TeX root = ../../packet-and-principal-models.tex
\section{Introduction}\label{sec:intro}
A rational CM Hodge structure does not determine an integral polarized
abelian variety, and neither determines its location on a specified curve.
This paper separates these three questions for one finite packet on
McMullen's curve. We first determine its intersection with the normalized
curve. We then construct principal models of the three surviving fibers
and distinguish their integral maps and polarizations.

The packet is fixed by the extension $M/L$, the ordered embeddings and the
norm generator below. All conjugacy counts use $\SL_2(\mathcal O_L)$.
Replacing this group by $\GL_2(\mathcal O_L)$ or forgetting a field marking
changes the classification problem. No assertion concerns every CM point
of the curve or every Hecke packet of level $43$.

\subsection{The three questions and their answers}
The finite packet has $3520$ trace-labelled matrix classes and $368$
quotient points. Its intersection with the normalized curve is exactly
$A,B,C$, subject to the execution premise stated below. The original
fibers have polarization type $(1,7,7,21,21,21)$, not principal type.
A common self-dual lattice gives isogenous principal models on a connected
finite cover. For the B/C principal models the minimum unrestricted
isogeny degree is $7$, whereas the minimum scalar-polarized degree is $64$.
These two minima measure different requirements on a map.

The table in Theorem~\ref{thm:A} presents all three geometric fibers
together. The main statements distinguish the finite intersection, the
Hodge consequences and the integral principal geometry. The supporting
Hodge note treats powers and mixed products; none of those extensions is
needed to prove the finite intersection.

\subsection{Fixed data and execution premise}
Write $z=\zeta_{42}$, $M=\mathbb Q(z)=\mathbb Q(\zeta_{21})$,
$L=M^+$ and $t=z+z^{-1}$. Order the real embeddings by
$1,5,11,13,17,19$ and put
\[
 \pi=t^4+t^3-4t^2-2t+5,\qquad \alpha_0=1+z+z^4,
 \qquad \alpha_k=z^k\alpha_0,\quad c_k=\alpha_k+\bar\alpha_k.
\]
Then $N_{M/L}(\alpha_0)=\pi$ and $N_{L/\mathbb Q}(\pi)=43$.
The symbols $A,B,C$ denote normalized vertex classes; $A_X$ denotes the
trace-polarized sixfold at $X$, and $B_X$ its specified principal model.
Let $F_t:\mathbb H\to\mathbb H^6$ be McMullen's modular embedding
with the component normalization of Supplement, Section~\externalref{supplement}{app:components}.
The normalized curve image is the image of $F_t$ in
$\SL_2(\mathcal O_L)\backslash\mathbb H^6$. Its quadrilateral
monodromy is the integral, orientation-preserving group used below.
The normalized generators are $X_t=h_tXh_t^{-1}$ with
$h_t=\operatorname{diag}(t,1)$. The fourth quadrilateral vertex is treated
in Proposition~\ref{prop:DBsame}; it is not a fourth quotient class.

\begin{executionhypothesis}\label{hyp:packetexecution}
The packet intersection claims use the retained exact search executions
identified in Supplement, Section~\externalref{supplement}{app:computer}. For the historical commands and
continuation cohort lacking complete end-of-execution guards, we assume
that the archived mathematical sources and runtime were not transiently
replaced during execution. The accepted source correspondence and final
population join are those of the final packet adoption identified there.
Later hashes and small control runs do not establish the missing historical
guards or constitute a fresh execution of the full suite.
This premise concerns the retained execution evidence; it is additional
to the finite-domain and exact-arithmetic arguments in the proofs.
\end{executionhypothesis}

\begin{theorem}[The fixed marked norm-$\pi$ packet]\label{thm:packetintro}
Under the packet execution hypothesis above, the following holds.
Fix the extension $M/L$, its $L$-action and ordered six real embeddings,
and the generator $\pi$ and elements
$\alpha_k=\zeta_{42}^k\alpha_0$ for $0\leq k<42$ specified above.
The integral matrices of determinant $\pi$ and trace
$c_k=\alpha_k+\overline{\alpha}_k$ have exactly $3520$
trace-labelled $\SL_2(\OL)$-conjugacy classes. Their upper fixed tuples
give exactly $368$ distinct points of
$\SL_2(\OL)\backslash\HH^6$.
In the normalized trace-polarized McMullen curve image, precisely three
of these points occur: $A,B,C$. There are $365$ exclusions, consisting
of thirteen maximal-multiplier points and all $352$ proper-multiplier
points. The complete $112$-point $d=7$ subpacket has three hits and
$109$ exclusions.
\end{theorem}
\begin{proof}
The arithmetic assertion is Theorem~\ref{thm:fullpacketarithmetic}.
The geometric assertion is Theorem~\ref{thm:fullpacketgeometry}, using
the retained exact certificates and execution-provenance assumptions
specified in Supplement, Section~\externalref{supplement}{app:computer}.
\end{proof}

\begin{theoremA}\label{thm:A}
Let $z_X$ be the normalized vertex $X\in\{A,B,C\}$ and set
$A_X=\C^6/(\OL z_X+\OL)$, with trace polarization
$E_1((a,b),(a',b'))=\Tr_{L/\Ql}(a'b-ab')$.
Each has polarization type $(1,7,7,21,21,21)$ and a compatible
$K=\Ql(\sqrt{-3})$ action of signature $(3,3)$ and Hermitian
discriminant $-1$. Their actual geometric CM types and Hodge algebras are
\[
\begin{array}{c|c|c}
X&\Phi_X&\operatorname{Hdg}^*(A_X)\\ \hline
A&\{1,5,17,19,29,31\}&\langle\text{divisors}\rangle\\
B&\{1,5,23,25,29,31\}&\langle\text{divisors},W_K(A_B)\rangle\\
C&\{1,5,11,13,17,19\}&\langle\text{divisors},W_K(A_C)\rangle
\end{array}
\]
The A type is induced and the B/C types are primitive.
Consequently the Hodge conjecture holds for all three fibers in every
codimension. Their common self-dual lattice modification supplies
principal models $B_A,B_B,B_C$ on one connected finite-cover RM family,
with degree-$453789$ isogenies preserving the rational CM Hodge
structures, the indicated $K$ signatures and discriminants.
\end{theoremA}

\noindent\emph{Proof locator for Theorem~\ref{thm:A}.}
The trace polarization is computed in Proposition~\ref{prop:E1type}.
The B calculation is in Sections~\ref{sec:actualcmtype}--\ref{sec:fullhodge};
Theorem~\ref{thm:Ahodge} and Proposition~\ref{prop:Cvertex} give A and C.
Proposition~\ref{prop:commonprincipal} and
Theorem~\ref{thm:principalfamily} supply the common principal family.

\begin{theorem}[Principal models and their integral distinctions]
\label{pp:thm:introgeometry}
The specified principal models $B_A,B_B,B_C$ are pairwise nonisomorphic
even as unpolarized complex abelian varieties. For maps $B_B\to B_C$,
the least isogeny degree is $7$. Among these maps pulling back the given
target principal polarization to a positive integer multiple of the
source polarization, the least multiplier is $2$, with exactly six
degree-$64$ maps sharing one kernel. Under the normalization of
Proposition~\ref{pp:prop:homideal}, their integral Hom module is
the target-$R$ ideal $\mathfrak a$ of class order two.

On each of the two underlying varieties $B_B,B_C$, exactly two
isomorphism classes of principal polarizations occur, represented by
the given polarization and its $4-t^2$ twist. These give four distinct
polarized isomorphism classes after forgetting all CM and RM markings.
The same twist extends over the principal family; on the full coarse
principal RM component it is the nontrivial involution
$w_i\mapsto w_i/\sigma_i(4-t^2)$. It preserves the entire principal
curve image. On each of $B_A,B_B,B_C$ with its prescribed RM action
there are exactly two compatible principal-polarization classes;
the resulting six Weil points all lie on this same image and remain
distinct after forgetting all CM and RM markings. For each of the four
B/C polarization pairs, the least scalar multiplier remains $2$,
with six degree-$64$ maps in the direction $B_B\to B_C$; their two possible kernels are complementary
order-$64$ subgroups of $B_B[2]$.
\end{theorem}
\begin{proof}
See Theorem~\ref{thm:principalfamily} and
Propositions~\ref{pp:prop:unpolarizeddistinct},
\ref{pp:prop:bcisogeny}, \ref{pp:prop:mindegree},
\ref{pp:prop:homideal}, \ref{pp:prop:actualpolarizations}
and~\ref{pp:prop:unitfamily}, together with
Propositions~\ref{pp:prop:unitcurve}, \ref{pp:prop:sixprincipal},
\ref{pp:prop:unmarkeda}, \ref{pp:prop:polarizationsquare} and
Corollary~\ref{pp:cor:weilimage}. Arbitrary fine-level lifts are not
part of the involution assertion.
\end{proof}

\begin{remark}[Principal-moduli normalization]\label{rem:WKnormalization}
The original lattice $\mathcal O_L^2$ with its trace form is not a
principally polarized model. The self-dual modification changes the lattice
through an isogeny; passage to a finite cover alone would not do so.
The resulting principal family lies in the component determined by that
self-dual module. Its identification with Mostaed's named principally
polarized locus remains open. The Hodge conclusions at the three special
fibers do not require that further identification, and no global
$K$-action on the family is asserted.
\end{remark}

\subsection{Proof strategy and relation to prior work}
The arithmetic classification includes every module overorder and
orientation, then identifies the two characteristic roots at the level of
quotient points. For geometry, two simultaneous target regions reduce curve
membership to finite exact searches. Both regions are needed: separate
choices of a reflected chamber in different coordinates are not allowed.
The positive vertex realizations and exhaustive empty searches prove the
intersection theorem.

The principal construction instead uses the discriminant module of the
trace lattice. A common nonfree maximal isotropic subgroup is stable under
the three specialized quadratic actions. Its stabilizer gives the finite
cover on which the lattice modification descends. Integral lattice and
unit calculations then distinguish the models, maps and twists.

The Hodge consequences use divisor classes and the split-Weil input of
Schoen and Markman \cite{Schoen1998,Markman2025}. OpenAI's general CM
manuscript \cite{OpenAICM2026}, if valid, implies the rational algebraicity
conclusions; it is not a proof input here. Goodson's factor comparison and
Saller's related discussion \cite{Goodson2024,Saller2026} also precede any
claim of novelty for the rational CM types. The contribution examined here
is the fixed packet and its integral geometry, not new abstract isogeny
classes. Precise historical priority remains a separate literature question.
The reported Aoki route in Saller is not used as a substitute proof.

The trace-dual and unit-norm description of polarizations is classical.
For comparison, Li--Mantovan--Pries prove uniqueness, when a principal
polarization exists, for simple maximal-order CM types under specified
unit-sign hypotheses \cite[Prop.~4.5]{LiMantovanPries2021}. Their
same-field example $y^7=x^3-1$ is a known principally polarized Jacobian
\cite[Prop.~6.11 and Lemma~6.12]{LiMantovanPries2021}.
Our B/C principal models have a proper endomorphism order. The results
here concern their exact integral classes, maps and realization on the
specified curve; neither the general polarization method nor same-field
principal existence is claimed as new.

The constructed principal family is not identified with Mostaed's named
principal normalization \cite{Mostaed2026}. It has a different integral
lattice. Corrections relevant to that distinction are recorded in the
supplement; they are secondary to the positive results of this paper.
